\chapter{Dataset}
\label{cap:dataset}

This chapter describes the dataset used to train and evaluate every classifier in the thesis.
The dataset is published on Hugging Face under \texttt{Marek324/speech-music-classification} as four configurations (\texttt{mini}, \texttt{mid}, \texttt{full}, \texttt{crit}) and is reproducible from public sources by running \texttt{scripts/dataset/build.py}.

\section{Audio-signal taxonomy}
\label{sec:dataset_overview}

Al-Shoshan~\cite{ALSHOSHAN200695} proposes the following categorization of audio signals:
\begin{enumerate}
    \item \textit{``Speech signal composed of a single talker within a specific time period.''}
    \item \textit{``Pure music signal without any speech component.''}
    \item \textit{``Mixture of single-talker speech and background music.''}
    \item \textit{``Songs, i.e., a mixture of music and a singer's voice.''}
    \item \textit{``Singing without musical accompaniment.''}
    \item \textit{``Abnormal music, including unconventional cadences, human whistles, or non-musical sounds embedded as part of a musical structure.''}
    \item \textit{``Speech signal composed of two or more speakers talking simultaneously.''}
    \item \textit{``Non-speech and non-music signals, such as car, motor, or fan sounds.''}
    \item \textit{``Complex sound mixtures involving multiple speakers, singers, and musical sources.''}
\end{enumerate}

We map these nine categories to three classifier output classes:
\begin{itemize}
    \item \textbf{Speech} (categories 1, 3, 7): voiced human speech as the salient signal, including in the presence of background music or other concurrent speakers.
    \item \textbf{Music} (categories 2, 4, 5, 6, 9): instrumental, vocal, acapella, or atypical-music. The line between category 4 (vocal music) and category 1 (speech with music) is resolved by labeling the dominant signal.
    \item \textbf{Inactive} (category 8 and silence): non-speech non-music audio. Named ``inactive'' rather than ``noise'' because the operationally-relevant property is that downstream pipelines (ASR, codec selector) should not be activated on this material.
\end{itemize}

The 3-class formulation differs from the 2-class formulation of the historical reference papers (Lavner-and-Ruinskiy, Khonglah-and-Prasanna, Lemaire-and-Holzapfel — all 2-class).
The deployment scenarios this thesis cares about — streaming gates in front of ASR or codec selection — need a separate inactive class so that silence and ambient noise can be routed correctly.

\section{Composition and Construction}
\label{sec:dataset_composition}

The full-tier dataset is built deterministically from seven source families.
The build script (\texttt{scripts/dataset/build.py}) reads pinned source identifiers from \texttt{scripts/dataset/sources.toml}; the train / validation / test split assignment uses a Bresenham-style interleaving generator that guarantees every subclass appears in every split regardless of source size.

\subsection*{Speech}

\textbf{LibriSpeech clean} (\texttt{train.360}) supplies the \texttt{speech\_clean} subclass; the \texttt{full} tier draws $960$ minutes ($16$\,h) from LibriSpeech's $\sim 360$\,h clean training split.
\textbf{LibriSpeech other} (\texttt{train.500}) supplies \texttt{speech\_dirty}; the \texttt{full} tier draws $480$ minutes.
The Silero VAD~\footnote{\url{https://github.com/snakers4/silero-vad}} runs over each clip to identify speech-active regions; clips are split at silence boundaries, and only the speech-active segments are retained.

\subsection*{Music}

\textbf{Free Music Archive (FMA)} contributes six genres — \texttt{music\_instrumental}, \texttt{music\_electronic}, \texttt{music\_pop}, \texttt{music\_rock}, \texttt{music\_hip-hop}, \texttt{music\_folk} — at $342.86$ minutes each, summing to $\sim 2\,057$ minutes.
We use the Hugging Face mirror \texttt{rpmon/fma-genre-classification} with genre filtering.
\textbf{bel\_canto} (\texttt{ccmusic-database/bel\_canto}) supplies \texttt{music\_acapella} at $342.86$ minutes; this is the music subclass closest to speech in feature space and is the canary for any speech-vs-vocal-music confusion.
All music sources are labeled by a custom \texttt{MusicLabeler} that combines pyannote-audio's segmentation pipeline (\texttt{pyannote/segmentation-3.0}) with an RMS-based loudness threshold (per-hop RMS over $20$\,ms windows, noise floor at the $20$th percentile $\times\,0.35$); a frame is labeled \texttt{music} if either pyannote returns \texttt{SPEECH} for it (instrumental music is similar enough to speech to register on this VAD) or its RMS exceeds the per-clip noise floor.

\subsection*{Inactive}

\textbf{DEMAND} (\texttt{voice-biomarkers/DEMAND-acoustic-noise}) is a corpus of $18$ environmental noise recordings.
The \texttt{full} tier draws $1\,200$ minutes from DEMAND, achieved by re-using the recordings with non-overlapping windows.
The DEMAND material is curated to contain only environmental noise, so labeling is trivial: every frame is tagged \texttt{inactive} via a no-op \texttt{SilenceLabeler} that trusts the source.

\subsection*{Synthetic augmentations}
\label{sec:dataset_augmentations}

Four synthetic subclasses fill in conditions that no good real-audio source covers.
Each is generated deterministically at build time by \texttt{scripts/dataset/augmentation.py}.
\textbf{\texttt{speech\_multispeaker}} ($180$ min) overlays two LibriSpeech clips from different speakers at relative gain $-5$ to $+5$\,dB (one speaker is up to $5$\,dB louder than the other); LibriMix-style.
\textbf{\texttt{speech\_som}} (speech over music, $180$ min) overlays a LibriSpeech clip on an FMA clip at a fixed $-10$\,dB music-relative-to-speech gain (music is $10$\,dB quieter than speech).
\textbf{\texttt{speech\_msom}} (multi-speaker over music, $120$ min) combines the two preceding constructions.
\textbf{\texttt{speech\_noisy}} ($480$ min) overlays a LibriSpeech clip on a DEMAND noise clip at SNR $5$ to $+20$\,dB.
All synthetic subclasses carry the speech class label.

\section{Tier strategy}
\label{sec:dataset_tier_strategy}

The dataset is built in three sizes plus the special crit configuration.
\textbf{\texttt{mini}} ($\sim 60$ min) is for CI and end-to-end smoke testing.
\textbf{\texttt{mid}} ($\sim 1\,200$ min, $20$\,h) is for ablation experiments where compute budget is the binding constraint; a single training run takes $\sim 20$ min on an A6000.
\textbf{\texttt{full}} ($\sim 6\,000$ min, $100$\,h) is the canonical training and evaluation surface; training takes $\sim 1$\,h on the A6000.
All chapter-6 results are on the full tier.

\subsection*{Bresenham split assignment}

The train / validation / test split uses a Bresenham-style interleaving generator (\texttt{scripts/dataset/process.py::\_bresenham\_split\_iter}) that distributes clips from each source proportionally to a configured target ratio (default $80\,\%/10\,\%/10\,\%$).
The interleaving is per-source per-subclass.
Earlier sequential-cutoff approaches assigned clips to splits based on cumulative-minutes thresholds and failed for sources that exhausted their available data before reaching the train cutoff (those sources contributed only to train).
At one development point the full-tier test split had only three of the fourteen subclasses populated.
The Bresenham generator guarantees every subclass appears in every split, regardless of source size.

\section{Statistics}
\label{sec:dataset_statistics}

\Cref{tab:dataset_full_tier_targets} reports the per-subclass design targets.

\begin{table}[h]
\centering
\begin{tabular}{lr}
\toprule
\textbf{Subclass} & \textbf{Target (min)} \\
\midrule
\texttt{speech\_clean}        & $960$ \\
\texttt{speech\_dirty}        & $480$ \\
\texttt{speech\_multispeaker} & $180$ (synthetic) \\
\texttt{speech\_som}          & $180$ (synthetic) \\
\texttt{speech\_msom}         & $120$ (synthetic) \\
\texttt{speech\_noisy}        & $480$ (synthetic) \\
\midrule
Each music subclass ($\times 7$) & $342.86$ \\
\midrule
\texttt{noise}                & $1\,200$ \\
\midrule
\textbf{Total}                & $\mathbf{6\,000}$ \\
\bottomrule
\end{tabular}
\caption{Per-subclass target minutes for the \texttt{full} tier. Class proportions: $40\,\%$ speech / $40\,\%$ music / $20\,\%$ inactive.}
\label{tab:dataset_full_tier_targets}
\end{table}

The realized composition (\cref{tab:dataset_class_summary}) tracks the design targets to within Silero VAD trim variance.

\begin{table}[h]
\centering
\begin{tabular}{lrr}
\toprule
\textbf{Class} & \textbf{Duration (min)} & \textbf{Share (\%)} \\
\midrule
\textbf{Speech (Total)} & $2\,401.09$ & $40.3\,\%$ \\
\textbf{Music (Total)} & $2\,362.67$ & $39.6\,\%$ \\
\textbf{Inactive (Total)} & $1\,200.02$ & $20.1\,\%$ \\
\midrule
\textbf{Total} & $\mathbf{5\,963.77}$ & $100.0\,\%$ \\
\bottomrule
\end{tabular}
\caption{Realized \texttt{full}-tier composition aggregated to class level (from \texttt{results/dataset\_stats.txt}). Per-subclass breakdown is in the appendix.}
\label{tab:dataset_class_summary}
\end{table}

The realized $40 / 40 / 20$ split closely matches the design $40\,\%$ speech / $40\,\%$ music / $20\,\%$ inactive target.

\section{Critical tier}
\label{sec:dataset_critical_tier}

The \texttt{crit} tier is a separately-built, hand-curated set of approximately $120$ clips ($\sim 30$ min total) drawn from sources outside the full-tier training distribution and chosen specifically to expose adversarial conditions.
It is constructed by \texttt{scripts/dataset/build.py --only-critical-set}, which is backed by a dedicated pipeline under \texttt{scripts/dataset/crit/}.

The crit set is manifest-driven.
\texttt{scripts/dataset/crit/manifest.toml} enumerates every clip with its source file, class, subclass, and one of two label sources: either an inline \texttt{labels = [\{label, start, end\}, \dots]} list of hand-curated frame-level annotations (used for non-trivial cases like \texttt{music\_beatbox}, where vocal-only ``music'' lines must be tracked frame-by-frame against intermittent spoken interjections), or \texttt{detector = "vad" | "music" | "silence"} which routes the clip through the same Silero / pyannote+RMS / no-op labelers used for the standard tiers when the source material is uniform enough.

Synthetic crit content is generated by dedicated scripts: \texttt{make\_noisy\_speech.py} mixes clean speech recordings with noise sources at fixed SNRs ($0$\,dB and $-3$\,dB) to produce \texttt{speech\_noisy\_0db} / \texttt{speech\_noisy\_neg3db}, while \texttt{make\_switching.py} and \texttt{make\_switching\_3class.py} generate switching clips with controlled label-change cadences ($1\,000$, $500$, $250$, and $125$\,ms) for the temporal-behavior evaluation in \cref{sec:eval_streaming}.

The crit tier covers \texttt{noise\_babble} (multi-speaker conversational noise), \texttt{noise\_machinery} (broadband mechanical), \texttt{noise\_outdoor} (traffic, urban hum), \texttt{noise\_wildlife} (birds, included as a control), \texttt{noise\_ambient} (room tone), \texttt{speech\_whisper} (voiceless speech), \texttt{speech\_noisy\_0db} and \texttt{speech\_noisy\_neg3db} (challenging-SNR speech), \texttt{speech\_msom} (multi-speaker over music), \texttt{music\_extreme} (jumpup, speedcore — extreme BPM and distortion that surface as broadband noise on log-mel), and \texttt{music\_beatbox} (vocal-only music with intermittent speech).

The crit tier is held out from training entirely.
No clip in it appears in the training, validation, or standard test splits of any tier.
The clip distribution is hand-selected rather than naturally weighted, so the crit-tier macro F1 is a stress test, not a deployment-scenario performance estimate.
The crit-tier evaluation methodology is in \cref{sec:eval_critical}; the chapter-6 results in \cref{sec:exp_critical_results}.
